% !TEX root = task-3-9-2.tex
% Editable task bank. Append new BankTask blocks at the end of the document.
\documentclass[12pt,a4paper]{article}
\usepackage{physics-textbook}
\usepackage{graphicx,needspace}
\usetikzlibrary{decorations.pathmorphing}
\newif\ifShowSolutions
\ifdefined\TutorEdition
  \ShowSolutionstrue
\else
  \ShowSolutionsfalse % Change to \ShowSolutionstrue to display solutions.
\fi
\newcounter{banktask}
\newsavebox{\BankTaskBox}
\newcommand{\BankTaskContent}[3]{%
  \noindent\textbf{\thebanktask.}\quad #2\par
  \ifShowSolutions
    \Needspace{50mm}% Keep the solution heading with its opening text or diagram.
    #3\par
    \ifdefined\TutorEdition
      \ifcsname TutorSource#1\endcsname
        \nopagebreak\smallskip{\small\csname TutorSource#1\endcsname\par}%
      \fi
    \fi
  \fi
}
\newcommand{\BankTask}[3]{%
  \refstepcounter{banktask}\par\medskip
  \sbox{\BankTaskBox}{\begin{minipage}{\linewidth}%
    \BankTaskContent{#1}{#2}{#3}\end{minipage}}%
  \ifdim\dimexpr\ht\BankTaskBox+\dp\BankTaskBox\relax>\dimexpr\textheight-20mm\relax
    \Needspace{35mm}\BankTaskContent{#1}{#2}{#3}%
  \else
    \Needspace{\dimexpr\ht\BankTaskBox+\dp\BankTaskBox\relax}%
    \noindent\usebox{\BankTaskBox}\par
  \fi
  \medskip
}
\begin{document}
\BookSubsection{3.9.2}{Сопротивление и удельное сопротивление проводников}


% Source sdamgia-7355: https://phys-ege.sdamgia.ru/problem?id=7355
\BankTask{sdamgia-7355}{
В школьной лаборатории есть два проводника круглого сечения. Удельное сопротивление первого проводника в $2$ раза больше удельного сопротивления второго проводника. Длина первого проводника в $2$ раза больше длины второго. При подключении этих проводников к одинаковым источникам постоянного напряжения за одинаковые интервалы времени во втором проводнике выделяется количество теплоты в $4$ раза большее, чем в первом. Каково отношение радиуса второго проводника к радиусу первого проводника? Источник тока считать идеальным.\par
}{
\smallskip\noindent\textbf{Решение}\par
Количество теплоты, выделяемое на проводнике, согласно закону Джоуля — Ленца: $Q= ( \frac{ U ^2 }{ R } ) \cdot t.$\par
Сопротивление проводника: $R= \frac{ \rho l}{ S } = \frac{ \rho l}{ \pi r ^2 } .$\par
Запишем соотношение: $\frac{ 1}{ 4 } = \frac{ Q_{1}}{ Q_{2} } = \frac{ R_{2}}{ R_{1} } = \frac{ \rho_{2} l_{2} r_{1} ^2 }{ r_{2} ^2 \rho_{1} l_{1} } = \frac{ 1}{ 2 \cdot 2 } \frac{ r_{1} ^2 }{ r_{2} ^2 } .$\par
Отсюда:$\frac{ r_{2}}{ r_{1} } =1.$\par
Ответ: $1$.\par
}

% Source sdamgia-10223: https://phys-ege.sdamgia.ru/problem?id=10223
\BankTask{sdamgia-10223}{
При подключении куска проволоки к полюсам батареи через нее течет ток силой $0{,}5 \,\text{А}$. Этот кусок проволоки сложили пополам, место сгиба разрезали. Затем разрезали каждый получившийся короткий провод на две равные части, зачистили концы и присоединили все эти части к полюсам батареи параллельно. Найдите силу тока, которая будет течь через батарею в этом случае. Внутреннее сопротивление батареи очень мало.\par
}{
\smallskip\noindent\textbf{Решение}\par
Сопротивление длинного однородного проводника обратно пропорционально площади его сечения: $R= \frac{ \rho l}{ S } .$ Всю проволоку сопротивлением $R$ разрезали на $4$ равные части сопротивлением $R/4$ каждая. Общее сопротивление такой цепи при параллельном соединении будет равно $R_{0}=R/16.$ Тогда по закону Ома в цепи будет течь ток, в $16$ раз больший первоначального $I=16I_{0}=8\,\text{А}.$\par
Ответ: $8$.\par
}

% Source sdamgia-24367: https://phys-ege.sdamgia.ru/problem?id=24367
\BankTask{sdamgia-24367}{
По проволочному резистору течет ток. Как изменятся при уменьшении длины проволоки в $4$ раза и увеличении силы тока вдвое тепловая мощность, выделяющаяся на резисторе, и его электрическое сопротивление?\par
$1$. Увеличится.\par
$2$. Уменьшится.\par
$3$. Не изменится.\par
\begin{center}\small\setlength{\tabcolsep}{3pt}\begin{tabular}{|p{0.45\linewidth}|p{0.45\linewidth}|}\hline
Тепловая мощность,
выделяющаяся на резисторе & Электрическое сопротивление
резистора\\\hline
 & \\\hline
\end{tabular}\end{center}\par
}{
\smallskip\noindent\textbf{Решение}\par
Сопротивление проводника зависит от его длины прямо пропорционально. Следовательно, при уменьшении длины в $4$ раза его сопротивление уменьшится в $4$ раза ($2$).\par
Тепловая мощность, выделяемая в проводнике, вычисляется по формуле $P=I^{2}R$. Сила тока увеличится в $2$ раза, а сопротивление уменьшится в $4$ раза, следовательно, тепловая мощность не изменится ($3$).\par
Ответ: $32$.\par
}

% Source sdamgia-45485: https://phys-ege.sdamgia.ru/problem?id=45485
\BankTask{sdamgia-45485}{
Отрезок провода с большим удельным сопротивлением подключён к клеммам источника постоянного напряжения. Отрезок провода заменили другим проводом такой же длины и из того же материала, но вдвое большего диаметра. Как изменились в результате такой замены сопротивление внешней цепи и сила тока во внешней цепи? Считать, что напряжение на внешней цепи остаётся неизменным.\par
$1$) увеличилась\par
$2$) уменьшилась\par
$3$) не изменилась\par
\begin{center}\small\setlength{\tabcolsep}{3pt}\begin{tabular}{|p{0.45\linewidth}|p{0.45\linewidth}|}\hline
Сопротивление & Сила тока\\\hline
 & \\\hline
\end{tabular}\end{center}\par
}{
\smallskip\noindent\textbf{Решение}\par
Электрическое сопротивление рассчитывается по формуле $R = \frac{ \rho l}{ S } .$ При увеличении диаметра провода площадь его сечения увеличивается, а сопротивление уменьшается ($2$).\par
Тогда сила тока в цепи увеличивается ($1$).\par
Ответ: $21$.\par
}

% Source sdamgia-11805: https://phys-ege.sdamgia.ru/problem?id=11805
\BankTask{sdamgia-11805}{
Для проведения лабораторной работы по обнаружению зависимости сопротивления проводника от его диаметра ученику выдали пять проводников, характеристики которых приведены в таблице. Какие два из предложенных ниже проводников необходимо взять ученику, чтобы провести данное исследование?\par
\begin{center}\small\setlength{\tabcolsep}{3pt}\begin{tabular}{|p{0.225\linewidth}|p{0.225\linewidth}|p{0.225\linewidth}|p{0.225\linewidth}|}\hline
№
проводника & Длина проводника & Диаметр проводника & Материал\\\hline
$1$ & $100 \,\text{см}$ & $1{,}0 \,\text{мм}$ & медь\\\hline
$2$ & $200 \,\text{см}$ & $0{,}5 \,\text{мм}$ & медь\\\hline
$3$ & $200 \,\text{см}$ & $1{,}0 \,\text{мм}$ & медь\\\hline
$4$ & $100 \,\text{см}$ & $0{,}5 \,\text{мм}$ & алюминий\\\hline
$5$ & $300 \,\text{см}$ & $1{,}0 \,\text{мм}$ & медь\\\hline
\end{tabular}\end{center}\par
}{
\smallskip\noindent\textbf{Решение}\par
Для определения зависимости сопротивления проводника от его диаметра необходимо выбрать проводники, у которых все остальные параметры (длина и материал) одинаковы. Это проводники с номерами $2$ и $3$.\par
Ответ: $23$.\par
}

% Source sdamgia-24955: https://phys-ege.sdamgia.ru/problem?id=24955
\BankTask{sdamgia-24955}{
На графике показана зависимость силы тока $I$, текущего в цилиндрическом медном проводнике с площадью поперечного сечения $3{,}4 \,\text{мм}^{2}$, от приложенного к его концам напряжения $U$. Удельное сопротивление меди $0{,}017 \,\text{Ом} \cdot \,\text{мм}^{2}/\,\text{м}$.\par
Чему равна длина этого проводника? (В метрах.)\par

\begin{center}\begin{tikzpicture}[book diagram,>=stealth]
\draw[task grid,xstep=1,ystep=.5](0,0)grid(3,3);\draw[task axis](0,0)--(3.4,0)node[right]{$U,\,\text{В}$};\draw[task axis](0,0)--(0,3.4)node[above]{$I,\,\text{А}$};\node[below left]at(0,0){$0$};\foreach\x/\s in{1/.5,2/1.0}{\node[below]at(\x,0){$\s$};}\foreach\y in{1,2}\node[left]at(0,\y){$\y$};\draw[task curve](.8,.8)--(2.75,2.75);
\end{tikzpicture}\end{center}
\par
}{
\smallskip\noindent\textbf{Решение}\par
Сопротивление проводника равно $R= \frac{ U}{ I } ,$ $R= \frac{ \rho l}{ S } .$ Следовательно:\par
\[l= \frac{ US}{ \rho I } = \frac{ 1 \cdot 3{,}4}{ 2 \cdot 0{,}017 } =100\,\text{м}.\]
Ответ: $100$.\par
}

% Source sdamgia-36761: https://phys-ege.sdamgia.ru/problem?id=36761
\BankTask{sdamgia-36761}{
Участок цепи состоит из двух последовательно соединенных цилиндрических проводников, сопротивление первого из которых равно $R$, а второго $2R$. Как изменится общее сопротивление этого участка цепи и тепловая мощность, выделяющаяся в первом проводнике, если удельное сопротивление и площадь поперечного сечения первого проводника увеличить вдвое, оставив без изменения его длину и напряжение на концах участка цепи? Сопротивление второго проводника не меняется.\par
$1$. Увеличится.\par
$2$. Уменьшится.\par
$3$. Не изменится.\par
\begin{center}\small\setlength{\tabcolsep}{3pt}\begin{tabular}{|p{0.45\linewidth}|p{0.45\linewidth}|}\hline
Общее сопротивление участка цепи & Тепловая мощность, выделяющаяся
в первом проводнике\\\hline
 & \\\hline
\end{tabular}\end{center}\par
}{
\smallskip\noindent\textbf{Решение}\par
А) При увеличении площади сечения и удельного сопротивления первого проводника в $2$ раза его сопротивление не изменится. Поскольку проводники соединены последовательно, то общее сопротивление участка цепи не изменится ($3$).\par
Б) Поскольку сопротивление первого проводника и напряжение не изменились, то выделяющаяся в нем тепловая мощность также не изменится ($3$).\par
Ответ: $33$.\par
}

% Source sdamgia-4206: https://phys-ege.sdamgia.ru/problem?id=4206
\BankTask{sdamgia-4206}{
Школьник изучал процесс протекания постоянного тока через проволоку постоянного поперечного сечения $2 \,\text{мм}^{2}$. Изменяя длину проволоки $L$, он измерял при помощи миллиомметра ее сопротивление $R$. Результаты его измерений приведены в таблице.\par
\begin{center}\small\setlength{\tabcolsep}{3pt}\begin{tabular}{|p{0.129\linewidth}|p{0.129\linewidth}|p{0.129\linewidth}|p{0.129\linewidth}|p{0.129\linewidth}|p{0.129\linewidth}|p{0.129\linewidth}|}\hline
$L$, см & $50$ & $70$ & $90$ & $110$ & $130$ & $150$\\\hline
$R$, мОм & $103$ & $140$ & $175$ & $228$ & $260$ & $298$\\\hline
\end{tabular}\end{center}\par
Пользуясь таблицей, определите удельное сопротивление металла, из которого была изготовлена проволока. (Ответ дать в омах квадратного миллиметра на метр, округлив до десятых.)\par
}{
\smallskip\noindent\textbf{Решение}\par
Сопротивление проводника длиной $L$ с площадью поперечного сечения $S$ и с удельным сопротивлением $\rho$ равно $R= \frac{ \rho L}{ S } .$ Для каждого измерения сосчитаем удельное сопротивление $\rho = \frac{ RS}{ L }$\par
\begin{center}\small\setlength{\tabcolsep}{3pt}\begin{tabular}{|p{0.129\linewidth}|p{0.129\linewidth}|p{0.129\linewidth}|p{0.129\linewidth}|p{0.129\linewidth}|p{0.129\linewidth}|p{0.129\linewidth}|}\hline
$L$, см & $50$ & $70$ & $90$ & $110$ & $130$ & $150$\\\hline
$R$, мОм & $103$ & $140$ & $175$ & $228$ & $260$ & $298$\\\hline
$\rho,\,\text{Ом} \cdot \,\text{мм} ^2 /\,\text{м}$ & $0{,}41$ & $0{,}40$ & $0{,}39$ & $0{,}41$ & $0{,}40$ & $0{,}40$\\\hline
\end{tabular}\end{center}\par
Усредним полученные выражения:\par
\[\frac{ 0{,}41 + 0{,}40 + 0{,}39 + 0{,}41 + 0{,}40 + 0{,}40}{ 6 } \approx 0{,}4\,\text{Ом} \cdot \,\text{мм} ^2 /\,\text{м}.\]
Ответ: $0{,}4$.\par
}

% Source sdamgia-38950: https://phys-ege.sdamgia.ru/problem?id=38950
\BankTask{sdamgia-38950}{
Нихромовый проводник длиной $l=l_{2}$ включён в цепь постоянного тока. К нему подключают вольтметр таким образом, что одна из клемм вольтметра всё время подключена к началу проводника, а вторая может перемещаться вдоль проводника. На рисунке приведена зависимость показаний вольтметра $U$ от расстояния $x$ до начала проводника. Как зависит от $x$ площадь поперечного сечения проводника? Ответ поясните, указав, какие физические закономерности вы использовали.\par

\begin{center}\begin{tikzpicture}[book diagram,>=stealth]
\draw[task axis](0,0)--(3.7,0)node[right]{$x$};\draw[task axis](0,0)--(0,3.1)node[above]{$U$};\node[below left]at(0,0){$0$};\draw[task projection](1,0)--(1,1.7)--(0,1.7);\draw[task projection](3,0)--(3,2.7)--(0,2.7);\node[below]at(1,0){$l_1$};\node[below]at(3,0){$l_2$};\node[left]at(0,1.7){$U_1$};\node[left]at(0,2.7){$U_2$};\draw[task curve](0,0)--(1,1.7)--(3,2.7);
\end{tikzpicture}\end{center}
\par
}{
\smallskip\noindent\textbf{Решение}\par
Так как проводник включен в цепь постоянного тока, то выполняется закон Ома для участка цепи $I= \frac{ U}{ R } ,$ где сила тока $I$ одинакова на всех участках цепи, напряжение на обоих участках растет, судя по графику, линейно. Сопротивление проводника равно $R= \frac{ \rho x}{ S } ,$ где $S$ — площадь поперечного сечения. В силу линейной зависимости напряжения от длины на каждом участке делаем вывод, что площадь поперечного сечения при $0 < x < l_{1}$ постоянно и равно $S_{1},$ при $l_{1} < x < l_{2}$ также постоянно и равно $S_{2}.$\par
Учитывая, что на первом участке напряжение растет быстрее, чем на втором, делаем вывод, что площадь поперечного сечения второго участка больше, чем у первого $S_{2} > S_{1}.$\par
Ответ: площадь поперечного сечения постоянна на каждом участке и $S_{2} > S_{1}.$\par
}

% Source sdamgia-41152: https://phys-ege.sdamgia.ru/problem?id=41152
\BankTask{sdamgia-41152}{
По цилиндрическому проводнику с площадью поперечного сечения $S=2 \,\text{мм}^{2}$ протекает электрический ток силой $I=1 \,\text{А}$. Параллельно проводнику подключён плоский конденсатор. Модуль напряжённости электрического поля между пластинами этого конденсатора равен $E=50 \,\text{В/м}$. Найдите удельное электрическое сопротивление $\rho$ данного проводника, если известно, что его длина в $n=1000$ раз больше расстояния между пластинами конденсатора.\par
}{
\smallskip\noindent\textbf{Решение}\par
$1$. По закону Ома для участка цепи напряжение между концами цилиндрического проводника равно $U_{R} = IR,$ где сопротивление проводника $R = \rho \frac{ l}{ S }$ ($l$ — длина проводника, $S$ — площадь его поперечного сечения).\par
$2$. Напряжение между пластинами плоского конденсатора $U_{C} = Ed,$ где $d$ — расстояние между пластинами.\par
$3$. Поскольку проводник и конденсатор соединены параллельно, $U_{C} = U_{R}.$ С учётом выражения для $R$ получим: $Ed = \frac{ I \rho l}{ S } .$ Отсюда\par
\[\rho = \frac{ ES}{ I } \cdot \frac{ d}{ l } = \frac{ ES}{ nI } = \frac{ 50\,\text{В/м} \cdot 2 \cdot 10 ^{ - 6 } \,\text{м} ^2 }{ 1000 \cdot 1\,\text{А} } = 0{,}1 \frac{ \,\text{Ом} \cdot \,\text{мм} ^2 }{ \,\text{м} } .\]
Ответ: $\rho = \frac{ ES}{ nI } = 0{,}1 \frac{ \,\text{Ом} \cdot \,\text{мм} ^2 }{ \,\text{м} } .$\par
}

% Source sdamgia-8557: https://phys-ege.sdamgia.ru/problem?id=8557
\BankTask{sdamgia-8557}{
Вольт-амперная характеристика лампы накаливания изображена на рисунке. При напряжении источника $12 \,\text{В}$ температура нити лампы равна $3100 \,\text{К}$. Сопротивление нити прямо пропорционально ее температуре. Какова температура нити накала при напряжении источника $6 \,\text{В}$?\par

\begin{center}\begin{tikzpicture}[book diagram,>=stealth]
\draw[task grid,xstep=.5,ystep=.4](0,0)grid(4,4.8);\draw[task axis](0,0)--(4.4,0)node[right]{$U,\,\text{В}$};\draw[task axis](0,0)--(0,5.2)node[above]{$I,\,\text{А}$};\foreach\x in{0,2,4,6,8,10,12,14}\node[below]at(\x/4,0){$\x$};\foreach\y in{.2,.4,.6,.8,1,1.2,1.4,1.6,1.8,2}\node[left]at(0,2*\y){$\y$};\draw[task curve]plot[smooth]coordinates{(0,0)(.5,1.3)(1,2.2)(1.5,2.8)(2,3.3)(2.5,3.7)(3,4)(3.5,4.25)(4,4.45)};
\end{tikzpicture}\end{center}
\par
}{
\smallskip\noindent\textbf{Решение}\par
$1$. При напряжении источника $U_{1}=12 \,\text{В}$ сила тока через лампу определяется из графика: $I_{1}=2 \,\text{А}$.\par
$2$. Сопротивление нити накала при этом определяется законом Ома: $R_{1}= \frac{ U_{1}}{ I_{1} } =6\,\text{Ом}$.\par
$3$. При напряжении $U_{2}=6 \,\text{В}$ сила тока равна $I_{2}=1{,}4 \,\text{А}$ (см. вольт-амперную характеристику).\par
$4$. Сопротивление нити накала при этом напряжении $R_{2}= \frac{ U_{2}}{ I_{2} } \approx4{,}3\,\text{Ом}$.\par
$5$. Поскольку сопротивление нити пропорционально температуре $R= \beta T,$ то $\frac{ T_{2}}{ T_{1} } = \frac{ R_{2}}{ R_{1} }$ и\par
$T_{2}=T_{1} \frac{ R_{2}}{ R_{1} } =T_{1} \frac{ U_{2}I_{1}}{ I_{2}U_{1} } =3100 \cdot \frac{ 6 \cdot 2}{ 1{,}4 \cdot 12 } \approx2214\,\text{К}$.\par
Ответ: $2214\,\text{К}$.\par
}

\end{document}
